\section{Implémentation et commentaire du code}
\label{sec:implementation}

Le code est organisé en quatre couches, avec une règle stricte : les couches
\code{core} et \code{coreset} ne connaissent pas Spark (pur Scala,
sérialisable, testable sans JVM Spark) ; toute la logique Spark vit dans
\code{streamkm.spark}, sous forme d'orchestration pure : aucun calcul
algorithmique n'y est dupliqué.

\begin{center}
\begin{tabular}{@{}ll@{}}
\toprule
Package & Contenu \\
\midrule
\code{streamkm.core}    & \code{Point}, \code{Distance}, \code{KMeans} (seeding, Lloyd, coût) \\
\code{streamkm.coreset} & \code{CoresetTree} (coreset tree), \code{BucketManager} (merge-and-reduce) \\
\code{streamkm.spark}   & \code{StreamKMPlusPlus} (portage Spark), \code{CostEvaluator} \\
\code{streamkm.experiments} & \code{QualityMetrics} (harnais d'expériences, section~\ref{sec:experiments}) \\
\bottomrule
\end{tabular}
\end{center}

\subsection{Structure de données : le point pondéré}

Toute la chaîne repose sur une seule structure de données, \code{Point},
volontairement minimale : un tableau de coordonnées et un poids (1.0 pour un
point brut du flux, la masse représentée pour un point de coreset).

\lstinputlisting[
  language=Scala,
  firstline=20, lastline=27,
  caption={\code{Point} : structure de données centrale (\code{streamkm/core/Point.scala})}
]{../StreamKMpp/src/main/scala/streamkm/core/Point.scala}

\subsection{Seeding $D^2$ pondéré}

\lstinputlisting[
  language=Scala,
  firstline=65, lastline=96,
  caption={\code{KMeans.seedPlusPlus} : seeding k-means++ pondéré}
]{../StreamKMpp/src/main/scala/streamkm/core/KMeans.scala}

\paragraph{Commentaire.} La ligne 73 tire le premier centre
\emph{proportionnellement au poids} plutôt qu'uniformément. Sur des points
bruts (poids uniformément égal à 1), les deux tirages coïncident exactement ;
mais dans le pipeline complet de StreamKM++, le seeding $D^2$ n'est
\emph{jamais} appliqué à des points bruts : il s'applique toujours au
coreset final, dont les poids sont très hétérogènes (la masse représentée
par chaque point). Un tirage uniforme y biaiserait l'initialisation vers les
régions les moins denses du coreset. Les lignes 77 à 79 pré-calculent la
distance de chaque point au premier centre (\code{closest}), mise à jour
incrémentalement à chaque nouveau centre (lignes 87 à 92) : c'est ce qui
donne la complexité $O(n\cdot k\cdot d)$ annoncée plutôt que
$O(n\cdot k^2\cdot d)$.

\subsection{Lloyd pondéré et son invariant de monotonie}

\lstinputlisting[
  language=Scala,
  firstline=141, lastline=178,
  caption={\code{KMeans.lloyd} : boucle principale (extrait), pondération et invariant}
]{../StreamKMpp/src/main/scala/streamkm/core/KMeans.scala}

\paragraph{Commentaire.} La ligne 149 accumule le coût \emph{avant} mise à
jour des centres (\code{currCost}), et non après : c'est cette quantité
précise dont la décroissance prouve la convergence de Lloyd (l'étape
d'affectation ne peut que diminuer le coût pour des centres fixés, et l'étape
de mise à jour ne peut que le diminuer pour une affectation fixée).
L'assertion de la ligne 173 rend cet invariant \emph{exécutable} : toute
régression future qui casserait la pondération (par exemple un retour
accidentel à une moyenne non pondérée) ferait échouer cette assertion en
mode debug plutôt que de produire silencieusement un résultat faux.

\subsection{Coreset tree : la descente pondérée par le coût}

\lstinputlisting[
  language=Scala,
  firstline=114, lastline=126,
  caption={\code{CoresetTree.descend} : choix de la feuille à scinder}
]{../StreamKMpp/src/main/scala/streamkm/coreset/CoresetTree.scala}

\paragraph{Commentaire.} La descente (ligne 121) compare un tirage aléatoire
au \emph{coût cumulé} du sous-arbre gauche (\code{node.left.cost}), jamais à
sa masse. Le cas \code{total > 0.0} faux (ligne 120) est une garde
défensive : si le coût du nœud est nul, tous ses points sont confondus avec
leur représentant et aucun choix n'a d'importance ; ce cas se produit
réellement sur des données très redondantes (voir le test \emph{"gère un
ensemble de points tous identiques"} de \code{CoresetSpec}).

\lstinputlisting[
  language=Scala,
  firstline=54, lastline=64,
  caption={\code{CoresetTree.build} : amorçage (extrait)}
]{../StreamKMpp/src/main/scala/streamkm/coreset/CoresetTree.scala}

La ligne 56 est la garantie structurelle centrale de toute la structure :
si l'ensemble d'entrée tient déjà dans le budget \code{m}, il \emph{est} son
propre coreset. C'est cette invariance de la masse totale que testent en
priorité les tests de la couche \code{coreset} (section~\ref{sec:tests}).

\subsection{Merge-and-reduce : consommation du flux par cascade de buckets}

\lstinputlisting[
  language=Scala,
  firstline=52, lastline=91,
  caption={\code{BucketManager.insert} : algorithme 2.5 (merge-and-reduce)}
]{../StreamKMpp/src/main/scala/streamkm/coreset/BucketManager.scala}

\paragraph{Commentaire.} Le déroulé en trois cas (ligne 62 à 90) suit
exactement la prose de la thèse (section 2.3.4) : si le bucket $B_1$ est
vide (ligne 64), on s'y contente de déplacer $B_0$ sans aucune réduction ;
$B_1$ contient alors $m$ points \emph{bruts} qui représentent bien $2^0\cdot
m$ points du flux, sans perte d'information. Si $B_1$ est plein (ligne 68),
la boucle \code{while} des lignes 76 à 87 implémente la cascade de
fusions-réductions : à chaque itération, on reconstruit un coreset de taille
$m$ (ligne 77, \code{CoresetTree.build}) à partir de l'union du niveau
courant et du niveau précédent, et on ne s'arrête que lorsqu'un niveau libre
est trouvé (ligne 79). L'assertion de la ligne 89
(\code{checkInvariants}), active uniquement avec le flag JVM \code{-ea},
vérifie après chaque insertion que la structure de buckets respecte bien
ses invariants de taille.

\subsection{Fusion de deux \code{BucketManager} : le \code{combOp} du portage Spark}

\lstinputlisting[
  language=Scala,
  firstline=154, lastline=166,
  caption={\code{BucketManager.merge} : fusion de deux structures indépendantes}
]{../StreamKMpp/src/main/scala/streamkm/coreset/BucketManager.scala}

\paragraph{Commentaire.} La ligne 162 utilise le générateur aléatoire de
\code{out}, l'objet construit avec la graine combinée \code{seed \^{}
other.seed} (ligne 157), et non celui de \code{this} ni de \code{other} :
ainsi, à chaque niveau de \code{treeReduce}, la construction du coreset
fusionné combine réellement la graine des deux opérandes de ce niveau
précis, plutôt que de ne dépendre que d'un seul d'entre eux. Cette propriété
est nécessaire pour que l'aléa reste indépendant à tous les niveaux d'une
fusion distribuée, et pas seulement au premier.

\subsection{Portage Spark : pourquoi \code{treeReduce} plutôt que \code{treeAggregate}}

\lstinputlisting[
  language=Scala,
  firstline=86, lastline=93,
  caption={\code{StreamKMPlusPlus.mergeCoresets} : fusion distribuée en arbre}
]{../StreamKMpp/src/main/scala/streamkm/spark/StreamKMPlusPlus.scala}

\paragraph{Commentaire.} Une alternative plus directe consisterait à
utiliser \code{RDD.treeAggregate(zeroValue)(seqOp, combOp, depth)} avec un
\code{zeroValue} unique partagé entre toutes les partitions. Nous nous en
écartons délibérément : \code{treeAggregate} clone ce \code{zeroValue} pour
chaque tâche, mais avec le \emph{même} état interne du générateur aléatoire
pour toutes les partitions, ce qui casse l'indépendance de graine par
partition nécessaire à l'expérience de scalabilité qui fait varier le
nombre de partitions à données fixes (sans que ce changement ne se confonde
avec un simple changement de graine). Nous construisons donc explicitement
un \code{BucketManager} par partition avec une graine dérivée de son indice
(\code{p.seed + idx}, ligne 88), puis les fusionnons par \code{treeReduce}
(ligne 92), qui offre la même structure de réduction en arbre bornée en
profondeur, sans ce défaut.

\subsection{Instrument de mesure : \code{CostEvaluator}}

\lstinputlisting[
  language=Scala,
  firstline=43, lastline=48,
  caption={\code{CostEvaluator.cost} : coût distribué, sans logique dupliquée}
]{../StreamKMpp/src/main/scala/streamkm/spark/CostEvaluator.scala}

\paragraph{Commentaire.} Cette fonction ne réimplémente aucune formule : elle
délègue entièrement à \code{KMeans.cost} (couche \code{core}, déjà testée),
appelée une fois par partition sur les centres reçus par diffusion
(\emph{broadcast}, implicite via la fermeture Scala capturant
\code{bcCenters}), puis somme les résultats partiels par un
\code{treeReduce}. Contrairement à la fusion de coresets, cette agrégation
est une simple somme, associative et exacte à la précision flottante près :
aucune composition d'erreur multiplicative n'intervient ici, à la différence
de section~\ref{sec:algorithms} et de l'analyse de tolérance en
section~\ref{sec:experiments}.

\subsection{Validation : suite de tests}
\label{sec:tests}

L'ensemble du code présenté ci-dessus est couvert par une suite de tests
ScalaTest exécutée via \code{sbt test} (code source complet en
Annexe~\ref{sec:appendix-code}) :

\begin{center}
\begin{tabular}{@{}llc@{}}
\toprule
Suite & Couverture & Tests \\
\midrule
\code{KMeansSpec}        & seeding, Lloyd, pondération, monotonie du coût & 7 \\
\code{CoresetSpec}       & coreset tree, merge-and-reduce, conservation de la masse & 14 \\
\code{StreamKMSpec}      & non-régression distribué/séquentiel, état multi-batch & 7 (1 ignoré\footnotemark) \\
\code{CostEvaluatorSpec} & coût distribué vs. séquentiel, invariance au partitionnement & 4 \\
\midrule
\textbf{Total} & & \textbf{32 (31 exécutés, 1 ignoré)} \\
\bottomrule
\end{tabular}
\end{center}
\footnotetext{Un test d'intégration bout-en-bout de \code{StreamKMPlusPlus.run}
via \code{MemoryStream} est marqué \code{ignore} sur la machine de
développement (Windows) : Structured Streaming y échoue lors de l'écriture
des métadonnées de checkpoint faute de \code{winutils.exe}/\code{HADOOP\_HOME},
une limitation d'environnement sans rapport avec la correction de
l'algorithme. Il est actif tel quel sous Linux/macOS ou dans un conteneur
Docker.}

Le test le plus significatif de \code{StreamKMSpec} compare le coût d'un
clustering obtenu en distribuant les données sur 8 partitions Spark à celui
d'une exécution séquentielle de référence sur les mêmes données, avec une
tolérance \emph{dérivée} (et non choisie arbitrairement) de la loi de
composition multiplicative des coresets présentée en
section~\ref{sec:solution} ; le détail du calcul est donné en
section~\ref{sec:experiments}.
